\section{绪论}

\subsection{研究背景与问题}

视觉语言模型的快速发展为多模态信息处理开辟了新的范式。CLIP 通过大规模对比预训练将图像和文本映射到统一的语义空间，实现了零样本分类和跨模态检索能力的突破。在此基础上，BLIP-2 引入 Q-Former 作为视觉-语言桥接器，LLaVA 通过简洁的线性投影实现了视觉指令跟随，Qwen-VL 系列则进一步优化了原生视觉-语言模型的端到端训练。然而，一个实际瓶颈逐渐显现：当前视觉语言模型的性能高度依赖输入图像的感知质量。在低分辨率、强模糊、高噪声和部分遮挡等视觉退化条件下，视觉编码器提取的特征发生畸变，语义理解和文本生成出现系统性偏差。

脑电图作为一种无创的神经成像技术，以毫秒级时间分辨率记录被试观看视觉刺激时的大脑皮层电活动。近年来的研究逐步揭示，EEG 信号中蕴含可解码的视觉语义信息——从物体类别到语义嵌入空间的映射，均已有实验验证。这自然引出一个问题：当视觉输入质量受损时，同步采集的 EEG 信号能否作为额外的信息通道，辅助视觉语言模型进行更准确的语义推理？

本课程设计正是在上述跨学科背景下展开，系统研究视觉退化条件下真实配对 EEG 信号对视觉语言模型语义识别和生成式描述的辅助作用。需要说明的是，本设计将 EEG 定位为视觉语义增强信号，其效果通过与 shuffled EEG（破坏时间同步关系但保留统计分布）和 random EEG（无神经语义信息的高斯噪声）的对照来客观评估。

\subsection{课程设计目标与意义}

本设计的目标是在 Thought2Text 小样本 EEG-image 数据集上，构建并评估一个视觉退化条件下的 EEG 辅助视觉语言理解系统。实验设计覆盖六种退化条件（clean、lowres16、strong\_blur、strong\_noise、occlusion50、mixed）和四种 EEG 控制模式（real EEG、vision-only、shuffled EEG、random EEG）。

在课程设计定位上：方法层面，本设计综合运用了深度学习模型构建（多分支 EEG 编码器、门控融合、交叉注意力）、预训练模型适配（CLIP 视觉/文本编码器、Qwen2-VL 生成模型）、参数高效微调（LoRA）和对比学习（InfoNCE）；系统层面，采用从受约束语义分类到 token 级融合再到生成式输出的递进设计，保证每个阶段都有可验证的中间结果；验证层面，通过 shuffled EEG 和 random EEG 的双重对照设计，保证性能差异可归因于真实 EEG-image 配对关系的存在与否。

\subsection{主要工作}

本设计完成了以下四项核心工作：

\begin{enumerate}
  \item 构建了包含六种退化条件和四种对照模式的 EEG 辅助 VLM 评估框架，采用 image-level split 消除训练-测试图像泄漏；
  \item 设计并实现了 A2 时间-频谱-空间 EEG 语义融合模型——通过三支路并行编码和后融合 Transformer 实现 EEG 多维度特征的端到端提取与交互；
  \item 提出了 VTF token 级融合方法，在 CLIP ViT visual token 层面通过交叉注意力注入 EEG 信息，生成增强视觉 token；
  \item 实现了基于 Qwen2-VL 的生成式视觉语言模型，通过 LoRA 高效微调和多候选重排序在单 GPU 上完成了小样本条件下的 free-form caption 生成。
\end{enumerate}

\subsection{报告结构}

本报告共分为十章。第 2 章梳理相关技术基础；第 3 章介绍数据集与任务设定；第 4 章阐述系统总体架构；第 5 章详述 A2 EEG 语义融合模型；第 6 章介绍 token 级 EEG 视觉增强方法；第 7 章描述生成式模型的方案设计与比较；第 8 章报告完整实验结果与分析；第 9 章概述工程实现；第 10 章总结全文并提出改进方向。
